% GENERATED FILE. Edit canonical lessons, then run npm run generate:books.
% canonical-source-hash: fnv1a64:e6b213a6b3ab7a9f
% canonical-lessons: AR-C363-nazariyya, AR-C363-fardiyya, AR-C363-mafhum, AR-C363-zahira, AR-C363-fassara

\chapter{Theories and concepts}
\label{ch:ar-theories-and-concepts}
\hlchaptermodality{363}

\begin{chapteropening}
\textbf{By the end of this chapter:} \emph{I can talk about a theory, a hypothesis, a concept and a phenomenon, and say what something means.}

Complete the last lesson of chapter 363: I can talk about a theory, a hypothesis, a concept and a phenomenon, and say what something means.
\end{chapteropening}

\section[\texorpdfstring{\ar{نظرية} (naẓariyya)}{نظرية (naẓariyya)}]{\ar{نظرية} (naẓariyya) — a theory}

\label{lesson:AR-C363-nazariyya}



\begin{quote}
Before the new one: say the Arabic for to deduce, then the Arabic for a contradiction.
\end{quote}

\subsection*{You'll want to know: \ar{نظرية}}
\textbf{\ar{نظرية}} — \emph{naẓariyya} — \textquotedblleft{}a theory\textquotedblright{}.

\begin{quote}
\ar{درسنا} \ar{نظرية} \ar{جديدة} \ar{في} \ar{الصف}.
\end{quote}

\begin{quote}
\emph{darasnā naẓariyya jadīda fī aṣ-ṣaff.}
\end{quote}

\begin{quote}
We studied a new theory in class.
\end{quote}

\ar{نظريا}, \emph{naẓariyyan}, in theory.

\begin{cousinweb}
\textbf{\ar{نظرية}} (\emph{naẓariyya}) is built on the native Arabic root \emph{n-ẓ-r}, \textquotedblleft{}to look, to regard\textquotedblright{} (Lane). \emph{Naẓariyya} is a modern abstract noun from \emph{naẓar}, looking, consideration: a way of seeing a thing, a theory.
\end{cousinweb}

\begin{grammarlens}[title={What you've built}]
One more word for giving reasons and opinions.
\end{grammarlens}

\subsection*{Your turn}
\begin{itemize}
\raggedright
  \item \emph{Say it:} \emph{naẓariyya}
  \item \emph{Say it:} \emph{naẓariyya}, once more
  \item \emph{Say it:} \emph{darasnā naẓariyya jadīda fī aṣ-ṣaff.} — We studied a new theory in class.
  \item \emph{Recall:} say the Arabic for to deduce, then the Arabic for a contradiction, then say \emph{naẓariyya} again
\end{itemize}

\begin{tcolorbox}[breakable,colback=teal!4,colframe=teal!35!black,title={Before you move on}]
What does \ar{نظرية} mean? (\textquotedblleft{}A theory\textquotedblright{}.) Say it once more.
\end{tcolorbox}
\section[\texorpdfstring{\ar{فرضية} (farḍiyya)}{فرضية (farḍiyya)}]{\ar{فرضية} (farḍiyya) — a hypothesis}

\label{lesson:AR-C363-fardiyya}



\begin{quote}
Before the new one: say the Arabic for a contradiction, then the Arabic for a theory.
\end{quote}

\subsection*{You'll want to know: \ar{فرضية}}
\textbf{\ar{فرضية}} — \emph{farḍiyya} — \textquotedblleft{}a hypothesis\textquotedblright{}.

\begin{quote}
\ar{أراد} \ar{الطالب} \ar{أن} \ar{يختبر} \ar{فرضيته}.
\end{quote}

\begin{quote}
\emph{arāda aṭ-ṭālib an yakhtabira farḍiyyatah.}
\end{quote}

\begin{quote}
The student wanted to test his hypothesis.
\end{quote}

\begin{cousinweb}
\textbf{\ar{فرضية}} (\emph{farḍiyya}) is built on the native Arabic root \emph{f-r-ḍ}, \textquotedblleft{}to make a notch; to make binding\textquotedblright{} (Lane). \emph{Farḍiyya} is a modern abstract noun from \emph{farḍ}, a supposition: an idea put forward to be tested.
\end{cousinweb}

\begin{grammarlens}[title={What you've built}]
One more word for giving reasons and opinions.
\end{grammarlens}

\subsection*{Your turn}
\begin{itemize}
\raggedright
  \item \emph{Say it:} \emph{farḍiyya}
  \item \emph{Say it:} \emph{farḍiyya}, once more
  \item \emph{Say it:} \emph{arāda aṭ-ṭālib an yakhtabira farḍiyyatah.} — The student wanted to test his hypothesis.
  \item \emph{Recall:} say the Arabic for a contradiction, then the Arabic for a theory, then say \emph{farḍiyya} again
\end{itemize}

\begin{tcolorbox}[breakable,colback=teal!4,colframe=teal!35!black,title={Before you move on}]
What does \ar{فرضية} mean? (\textquotedblleft{}A hypothesis\textquotedblright{}.) Say it once more.
\end{tcolorbox}
\section[\texorpdfstring{\ar{مفهوم} (mafhūm)}{مفهوم (mafhūm)}]{\ar{مفهوم} (mafhūm) — a concept}

\label{lesson:AR-C363-mafhum}



\begin{quote}
Before the new one: say the Arabic for a theory, then the Arabic for a hypothesis.
\end{quote}

\subsection*{You'll want to know: \ar{مفهوم}}
\textbf{\ar{مفهوم}} — \emph{mafhūm} — \textquotedblleft{}a concept\textquotedblright{}.

\begin{quote}
\ar{مفهوم} \ar{الصداقة} \ar{مختلف} \ar{عند} \ar{كل} \ar{شخص}.
\end{quote}

\begin{quote}
\emph{mafhūm aṣ-ṣadāqa mukhtalif ʿinda kulli shakhṣ.}
\end{quote}

\begin{quote}
The concept of friendship differs from person to person.
\end{quote}

Also an adjective, understood, clear: \ar{هذا} \ar{مفهوم}.

\begin{cousinweb}
\textbf{\ar{مفهوم}} (\emph{mafhūm}) is built on the native Arabic root \emph{f-h-m}, \textquotedblleft{}to understand\textquotedblright{} (Lane). \emph{Mafhūm} is the passive participle of \emph{fahima}, \textquotedblleft{}to understand\textquotedblright{}: what is understood, a concept.
\end{cousinweb}

\begin{grammarlens}[title={What you've built}]
One more word for giving reasons and opinions.
\end{grammarlens}

\subsection*{Your turn}
\begin{itemize}
\raggedright
  \item \emph{Say it:} \emph{mafhūm}
  \item \emph{Say it:} \emph{mafhūm}, once more
  \item \emph{Say it:} \emph{mafhūm aṣ-ṣadāqa mukhtalif ʿinda kulli shakhṣ.} — The concept of friendship differs from person to person.
  \item \emph{Recall:} say the Arabic for a theory, then the Arabic for a hypothesis, then say \emph{mafhūm} again
\end{itemize}

\begin{tcolorbox}[breakable,colback=teal!4,colframe=teal!35!black,title={Before you move on}]
What does \ar{مفهوم} mean? (\textquotedblleft{}A concept\textquotedblright{}.) Say it once more.
\end{tcolorbox}
\section[\texorpdfstring{\ar{ظاهرة} (ẓāhira)}{ظاهرة (ẓāhira)}]{\ar{ظاهرة} (ẓāhira) — a phenomenon}

\label{lesson:AR-C363-zahira}



\begin{quote}
Before the new one: say the Arabic for a hypothesis, then the Arabic for a concept.
\end{quote}

\subsection*{You'll want to know: \ar{ظاهرة}}
\textbf{\ar{ظاهرة}} — \emph{ẓāhira} — \textquotedblleft{}a phenomenon\textquotedblright{}.

\begin{quote}
\ar{المطر} \ar{في} \ar{الصيف} \ar{ظاهرة} \ar{نادرة} \ar{هنا}.
\end{quote}

\begin{quote}
\emph{al-maṭar fī aṣ-ṣayf ẓāhira nādira hunā.}
\end{quote}

\begin{quote}
Rain in summer is a rare phenomenon here.
\end{quote}

Plural: \ar{ظواهر}, \emph{ẓawāhir}.

\begin{cousinweb}
\textbf{\ar{ظاهرة}} (\emph{ẓāhira}) is built on the native Arabic root \emph{ẓ-h-r}, \textquotedblleft{}to appear, to come out\textquotedblright{} (Lane). \emph{Ẓāhira} is the feminine active participle, what appears: a phenomenon.
\end{cousinweb}

\begin{grammarlens}[title={What you've built}]
One more word for giving reasons and opinions.
\end{grammarlens}

\subsection*{Your turn}
\begin{itemize}
\raggedright
  \item \emph{Say it:} \emph{ẓāhira}
  \item \emph{Say it:} \emph{ẓāhira}, once more
  \item \emph{Say it:} \emph{al-maṭar fī aṣ-ṣayf ẓāhira nādira hunā.} — Rain in summer is a rare phenomenon here.
  \item \emph{Recall:} say the Arabic for a hypothesis, then the Arabic for a concept, then say \emph{ẓāhira} again
\end{itemize}

\begin{tcolorbox}[breakable,colback=teal!4,colframe=teal!35!black,title={Before you move on}]
What does \ar{ظاهرة} mean? (\textquotedblleft{}A phenomenon\textquotedblright{}.) Say it once more.
\end{tcolorbox}
\section[\texorpdfstring{\ar{فسر} (fassara)}{فسر (fassara)}]{\ar{فسر} (fassara) — to interpret}

\label{lesson:AR-C363-fassara}



\begin{quote}
Before the new one: say the Arabic for a concept, then the Arabic for a phenomenon.
\end{quote}

\subsection*{You'll want to know: \ar{فسر}}
\textbf{\ar{فسر}} — \emph{fassara} — \textquotedblleft{}to interpret\textquotedblright{}.

\begin{quote}
\ar{فسر} \ar{المعلم} \ar{القصيدة} \ar{للطلاب}.
\end{quote}

\begin{quote}
\emph{fassara al-muʿallim al-qaṣīda li-ṭ-ṭullāb.}
\end{quote}

\begin{quote}
The teacher interpreted the poem for the students.
\end{quote}

Present: \ar{يفسر}, \emph{yufassir}.

The noun \ar{تفسير}, \emph{tafsīr}, an interpretation.

\begin{cousinweb}
\textbf{\ar{فسر}} (\emph{fassara}) is built on the native Arabic root \emph{f-s-r}, \textquotedblleft{}to uncover, to make plain\textquotedblright{} (Lane). It is Form II, to make plain: to explain what a thing means.
\end{cousinweb}

\begin{grammarlens}[title={What you've built}]
One more word for giving reasons and opinions.
\end{grammarlens}

\subsection*{Your turn}
\begin{itemize}
\raggedright
  \item \emph{Say it:} \emph{fassara}
  \item \emph{Say it:} \emph{fassara}, once more
  \item \emph{Say it:} \emph{fassara al-muʿallim al-qaṣīda li-ṭ-ṭullāb.} — The teacher interpreted the poem for the students.
  \item \emph{Recall:} say the Arabic for a concept, then the Arabic for a phenomenon, then say \emph{fassara} again
\end{itemize}

\begin{tcolorbox}[breakable,colback=teal!4,colframe=teal!35!black,title={Before you move on}]
What does \ar{فسر} mean? (\textquotedblleft{}To interpret\textquotedblright{}.) Say it once more.
\end{tcolorbox}
